% Generated draft: requires mathematical and source review.
\documentclass[11pt]{article}
\usepackage[margin=25mm]{geometry}
\usepackage{fontspec}
\setmainfont{Latin Modern Roman}
\usepackage{amsmath,amssymb,mathtools,unicode-math}
\setmathfont{Latin Modern Math}
\usepackage{enumitem,needspace,xurl,hyperref}
\setlength{\parindent}{0pt}
\setlength{\parskip}{5pt}
\setlength{\emergencystretch}{3em}
\newcommand{\N}{\mathbb{N}}
\newcommand{\Z}{\mathbb{Z}}
\newcommand{\Q}{\mathbb{Q}}
\newcommand{\R}{\mathbb{R}}
\newcommand{\C}{\mathbb{C}}
\newcommand{\F}{\mathbb{F}}
\newcommand{\A}{\mathbb{A}}
\newcommand{\PP}{\mathbb{P}}
\newcommand{\E}{\mathbb{E}}
\newcommand{\eps}{\varepsilon}
\newcommand{\dd}{\,\mathrm{d}}
\begin{document}
\section*{UnsolvedMath ID: 11000157}
\textbf{Source problem code:} catalogue entry
\subsection{Definitions and mathematical statement}
For a group $G$, its lower central series is $\gamma_1G=G$ and $\gamma_{r+1}G=[G,\gamma_rG]$. Its associated graded Lie ring is
\[
 \operatorname{gr}(G)=\bigoplus_{r\geq1}\gamma_rG/\gamma_{r+1}G,
\]
with bracket induced by group commutators.

Problem 1.1 is a research program rather than a fully quantified conjecture. The source says the preceding ideas concern the beginning of the lower central series of $I_g$, points out the lower central series of $K_g$, and proposes further study of the relation between the group structure of $M_g$ and $3$-manifold topology, with the possibility of new $3$-manifold invariants.

Here $M_g$ is described as a group, not a moduli space. The usual interpretation is a mapping class group, with $I_g$ its Torelli subgroup and $K_g$ its Johnson kernel, but none of these definitions or the surface conventions are supplied in the excerpt. The lower central series of either group must not be identified without proof with the Johnson filtration. No particular term, coefficient ring, computation, correspondence, class of $3$-manifolds, or invariance criterion is prescribed. The preceding ideas and existing subgroup/topology correspondence are also absent. Consequently the record supports this stated direction of investigation, not a newly invented assertion that a specific graded quotient defines a $3$-manifold invariant.
\subsection{Short English statement}
Investigate the lower central series of the Torelli and Johnson-kernel-type groups and their relationship with $3$-manifold topology, seeking possible new invariants. The excerpt states a program without a definite target theorem.
\subsection{Sources}
\begin{itemize}
\item Local registry: \url{data/unsolvedmath/problems.json}, permanent UnsolvedMath ID $11000157$, exact source code \texttt{catalogue entry}.
\item Farb, \emph{Problems on Mapping Class Groups and Related Topics}, Problem 1.1, PDF page 151: \url{https://www.math.uchicago.edu/~farb/papers/mcgbook.pdf}. Supplied access date: 2026-07-29; extraction: pdf-font-structure; repeated local numbering disambiguated by page among 329 headings.
\item The preceding constructions and group definitions are not included. No additional authoritative citation or individual attribution is supplied, and the registry literature summary is truncated.
\end{itemize}
\end{document}
